\documentclass[12pt]{article}

\usepackage[T1]{fontenc}
\usepackage[utf8]{inputenc}
\usepackage{geometry}
\usepackage{setspace}
\usepackage{parskip}

\geometry{margin=1in}
\onehalfspacing

\title{The Last Remaining Asymmetry:\\Fairness, Competition, and the Future of Fertility}
\author{Kangbin Yim\thanks{Email: yim@sch.ac.kr (Prof., SCH University, Korea)}}
\date{May 2026}

\begin{document}

\maketitle

For decades, fertility decline has been explained primarily through economics. Housing costs are too high. Education is too expensive. Employment is unstable.
The opportunity cost of raising children continues to increase.

Governments around the world have responded with subsidies, tax incentives, childcare programs, and various forms of family support. Yet despite these efforts, fertility rates continue to decline across many highly developed societies.

Perhaps this suggests that the problem is not purely economic. Perhaps something deeper is occurring beneath the surface. Most discussions about fertility focus on cost. Few focus on the structure of civilization itself.
Yet throughout history, civilization has followed a remarkably consistent direction.

It has gradually reduced asymmetry. The distance between rulers and citizens has narrowed. The distinction between aristocrats and commoners has weakened. Educational opportunities have become more accessible. Political participation has become more inclusive. Economic participation has become less dependent upon inherited privilege. In many ways, the history of civilization can be understood as a long process of expanding fairness.

But fairness produces a consequence that is often overlooked. Fairness does not eliminate competition. It universalizes competition. In a highly stratified society, many individuals are excluded from competition before it even begins. In a more symmetrical society, more people enter the same arena. 

The result is not less competition, but broader competition. The modern world is therefore not merely a fairer world. It is also a more competitive world. This transformation is especially visible in the relationship between men and women. For most of human history, men and women occupied substantially different social roles. The reasons were not purely cultural. They were also functional.

Survival depended heavily upon physical strength, environmental risk, warfare, hunting, and labor-intensive production. Men and women participated in society through different pathways because society itself was organized around different requirements.

Whether these arrangements were fair or unfair is a separate question. The important point is that the competition itself was often separated. Men and women were frequently competing within different domains. Modern civilization has steadily dismantled these distinctions. Industrialization reduced the importance of physical strength. Mass education expanded access to knowledge. Digital technology reduced barriers to information. Professional opportunities expanded. Economic independence expanded. Political rights expanded.

As a result, men and women increasingly participate within the same systems of competition. They attend the same schools. Compete for the same university admissions. Compete for the same professional positions. Compete for the same promotions. Compete for the same leadership roles. Compete for the same definitions of success. The trajectory is clear. The modern world is not merely creating greater equality. It is creating a shared competitive arena. 

Today, artificial intelligence and robotics are accelerating this process. Machines have already replaced much of the physical labor that once favored male strength. Artificial intelligence is beginning to automate forms of cognitive labor that once created other forms of differentiation. Robots increasingly perform physically demanding work. Algorithms increasingly perform analytical work. As these technologies advance, biological sex becomes progressively less relevant to economic productivity and professional achievement.

The long-term effect may be unprecedented. For the first time in human history, men and women may find themselves competing within almost entirely symmetrical social structures. And it is precisely at this moment that a paradox emerges. As more asymmetries disappear, one asymmetry becomes increasingly visible. Pregnancy and childbirth. Pregnancy occurs within the female body. Childbirth occurs within the female body. No political reform can eliminate this reality. No economic policy can completely redistribute it. No legal framework can fully equalize it. It remains a biological fact.

Paradoxically, the more symmetrical society becomes, the more visible this remaining asymmetry appears. The future challenge of fertility may therefore be less about economics alone and more about how civilization responds to the last remaining asymmetry.

\end{document}
